%=======================================================================
%  LABORATORIO N.o 2 - APRENDIENDO A RECTIFICAR SENALES
%  Curso : Circuitos Electronicos  (2026-2)
%  UCSP  - Facultad de Ingenieria - Ing. Electronica y de Telecomunicaciones
%  Docente: Ebert San Roman Castillo
%
%  Version 1 - Transcripcion del documento
%              "Laboratorio_2_Propedeucio EEB.doc".
%
%  Esta version NO anade contenido: reproduce la guia original. Solo se
%  corrigieron erratas ortograficas evidentes del .doc ("se parecia" ->
%  "se aprecia", "oscilocopio" -> "osciloscopio", "pareciar" -> "apreciar",
%  "senal cuadra" -> "senal cuadrada") y se dejo que LaTeX numere las
%  figuras: el original repite el numero 10 y el numero 11 en dos figuras
%  distintas cada uno, y con \ref las referencias del texto quedan
%  correctas por si solas.
%
%  Las cajas notatec / seguridad / ejemplo quedan definidas en el preambulo
%  para la futura revision tecnica, pero en esta version no se usan.
%
%  Compilar:  xelatex Laboratorio_2.tex    (dos veces)
%=======================================================================
\documentclass[11pt,a4paper]{article}

%------------------------------------------------------------- idioma
% Compilar con XeLaTeX (no con pdflatex): Trebuchet MS es una fuente TrueType
% del sistema y solo fontspec, bajo XeLaTeX o LuaLaTeX, sabe cargarla.
% inputenc y fontenc se omiten a proposito: XeLaTeX trabaja en UTF-8 de forma
% nativa y esos paquetes dan error con este motor.
\usepackage{fontspec}
\setmainfont{Trebuchet MS}[
  Path            = C:/Windows/Fonts/,
  Extension       = .ttf,
  UprightFont     = trebuc,
  BoldFont        = trebucbd,
  ItalicFont      = trebucit,
  BoldItalicFont  = trebucbi,
  Ligatures       = TeX
]
\usepackage[spanish,es-nodecimaldot]{babel}
\usepackage{microtype}

%----------------------------------------------------------- geometria
\usepackage[left=2.4cm,right=2.4cm,top=3.1cm,bottom=2.6cm,headheight=30pt,
            footskip=1.2cm]{geometry}

%------------------------------------------------------------ graficos
\usepackage{graphicx}
\usepackage{float}
\usepackage{placeins}
\usepackage[font=small,labelfont=bf,labelsep=period]{caption}
\graphicspath{{imagenes/}}

%-------------------------------------------------------------- tablas
\usepackage{array}
\usepackage{tabularx}
\usepackage{booktabs}
\usepackage{multirow}

%----------------------------------------------------- matematicas / SI
\usepackage{amsmath,amssymb}
\usepackage{siunitx}
\sisetup{
  output-decimal-marker = {,},
  per-mode              = symbol,
  range-phrase          = { a },
  range-units           = single,
  detect-weight         = true,
  detect-family         = true
}

%--------------------------------------------------------------- listas
\usepackage{enumitem}
\setlist[itemize]{itemsep=3pt,topsep=4pt,parsep=0pt}
\setlist[enumerate]{itemsep=8pt,topsep=4pt,parsep=0pt}

%--------------------------------------------------------------- cajas
\usepackage[most]{tcolorbox}
\usepackage{pgffor}

\newtcolorbox{notatec}[1]{enhanced,breakable,
  colback=blue!3!white, colframe=blue!55!black,
  fonttitle=\bfseries\small, title={Nota técnica --- #1},
  boxrule=0.6pt, arc=2pt, left=6pt,right=6pt,top=4pt,bottom=4pt}

\newtcolorbox{seguridad}[1]{enhanced,breakable,
  colback=red!3!white, colframe=red!60!black,
  fonttitle=\bfseries\small, title={Seguridad --- #1},
  boxrule=0.6pt, arc=2pt, left=6pt,right=6pt,top=4pt,bottom=4pt}

\newtcolorbox{ejemplo}[1]{enhanced,breakable,
  colback=green!3!white, colframe=green!45!black,
  fonttitle=\bfseries\small, title={Ejemplo de diseño --- #1},
  boxrule=0.6pt, arc=2pt, left=6pt,right=6pt,top=4pt,bottom=4pt}

%------------------------------------------------------------ comandos
% Lineas punteadas para respuestas del alumno: \lineasresp{n}
\newcommand{\lineasresp}[1]{%
  \par\vspace{6pt}%
  \foreach \i in {1,...,#1}{\noindent\rule{\linewidth}{0.4pt}\par\vspace{11pt}}%
}

%---------------------------------------------------- encabezado / pie
\usepackage{fancyhdr}
\pagestyle{fancy}
\fancyhf{}
\fancyhead[L]{\footnotesize UCSP --- Facultad de Ingeniería\\
              Ing. Electrónica y de Telecomunicaciones}
\fancyhead[R]{\footnotesize 2026-2 Circuitos Electrónicos\\
              Ebert San Román Castillo}
\fancyfoot[L]{\footnotesize Laboratorio N.\textsuperscript{o} 2 --- Aprendiendo a rectificar señales}
\fancyfoot[R]{\footnotesize Página \thepage\ de \pageref{LastPage}}
\renewcommand{\headrulewidth}{0.4pt}
\renewcommand{\footrulewidth}{0.4pt}
\usepackage{lastpage}

%------------------------------------------------------- numeracion
\renewcommand{\thesection}{\Roman{section}}
\renewcommand{\thesubsection}{\arabic{subsection}}

%-------------------------------------------------------- hiperenlaces
\usepackage[hidelinks,pdfusetitle]{hyperref}
\hypersetup{
  pdftitle={Laboratorio 2 - Aprendiendo a rectificar senales sinusoidales},
  pdfsubject={Circuitos Electronicos - Guia de practicas},
  pdfauthor={Ebert San Roman Castillo},
  pdfkeywords={electronica, diodo, rectificador, condensador, osciloscopio, Electronics Explorer Board}
}

\begin{document}

%=======================================================================
%  CABECERA DE LA GUIA
%=======================================================================
% El documento es de 11 pt, asi que \normalsize son los 11 pt del Word original.
\begin{center}
    {\normalsize\bfseries SEGUNDA UNIDAD: APRENDIENDO A RECTIFICAR SEÑALES}
\end{center}

\vspace{0.55cm}

% Cabecera identica a la del Laboratorio 1 (ver los comentarios de aquel archivo).
\begingroup
\renewcommand{\tabularxcolumn}[1]{m{#1}}
\renewcommand{\arraystretch}{2.0}
\setlength{\tabcolsep}{12pt}
\noindent
\begin{tabularx}{\textwidth}{|>{\centering\arraybackslash}m{2.2cm}|>{\centering\arraybackslash}X|>{\centering\arraybackslash}m{3.0cm}|}
\hline
\rule{0pt}{2.1cm}{\fontsize{58}{58}\selectfont\bfseries 2}
    & \begin{minipage}[c]{\linewidth}\centering
          {\large\bfseries Guía de Prácticas}\\[7pt]
          {\Large\bfseries Aprendiendo a Rectificar Señales Sinusoidales}
      \end{minipage}
    & {\Large\bfseries Nota}\\
\hline
\multicolumn{3}{|l|}{\rule{0pt}{1.3cm}\large\textbf{Grupo:}~\rule{4.6cm}{0.4pt}%
    \hfill\large\textbf{Fecha:}~\rule{4.6cm}{0.4pt}\hspace{2mm}}\\
\hline
\multicolumn{3}{|l|}{\rule{0pt}{1.2cm}\large\textbf{Alumno(s):}~\rule{11.2cm}{0.4pt}\hspace{2mm}}\\
\multicolumn{3}{|l|}{\rule{0pt}{1.2cm}\phantom{\large\textbf{Alumno(s):}}~\rule{11.2cm}{0.4pt}\hspace{2mm}}\\
\hline
\end{tabularx}
\endgroup

\vspace{1.2cm}

%=======================================================================
\section{Objetivos}
%=======================================================================
\begin{itemize}
    \item Continuar familiarizándose con el sistema de medición \emph{Electronics Explorer Board}.
    \item Reconocer los diferentes dispositivos de instrumentación electrónica.
    \item Reconocer componentes electrónicos tales como los diodos, resistencia, potenciómetro
          y capacitores.
    \item Comprender el uso de los diodos como rectificadores.
    \item Comprender la aplicación de los condensadores en la rectificación de una señal.
\end{itemize}

%=======================================================================
\section{Contenido teórico}
%=======================================================================
\subsection{Las medidas eléctricas}

Las medidas eléctricas utilizan métodos, dispositivos y cálculos usados para medir cantidades
eléctricas. La medición de cantidades eléctricas abarca un amplio uso; generalmente utilizan
transductores, que permiten medir propiedades físicas como la temperatura, presión, flujo,
fuerza, y muchas otras al convertirlas en señales eléctricas, las cuales pueden ser medidas y
registradas.

En esta práctica utilizaremos 2 herramientas básicas de la instrumentación electrónica, que son:

\begin{enumerate}
    \item El osciloscopio.
    \item Generador de señales.
\end{enumerate}

El osciloscopio es un aparato de medida que permite medir los valores de voltaje de un circuito
en función del tiempo; es un instrumento indispensable en el campo de la electrónica. Posee
controles de la división vertical, amplitud de la señal, y de la división horizontal, que es la
escala temporal, que se aprecian en la figura~\ref{fig:scope}.

\begin{figure}[H]
    \centering
    \includegraphics[width=0.50\textwidth]{image01.png}
    \caption{Controles del osciloscopio.}
    \label{fig:scope}
\end{figure}

Un generador de señales, o de funciones o de formas de onda, es un dispositivo electrónico de
laboratorio que genera patrones de señales periódicas o no periódicas tanto analógicas como
digitales. Se emplea normalmente en el diseño, prueba y reparación de dispositivos electrónicos.
La figura~\ref{fig:wavegen} muestra la pantalla de un generador de señales, en la que se pueden
visualizar los diferentes tipos de señales que puede producir el instrumento, como señales
sinusoidales, rectangulares, triangulares, dientes de sierra y otras.

\begin{figure}[H]
    \centering
    \includegraphics[width=0.29\textwidth]{image02.png}
    \caption{Pantalla de un generador de señales.}
    \label{fig:wavegen}
\end{figure}

\FloatBarrier
\subsection{Componentes electrónicos pasivos}

\subsubsection{El diodo}

Los semiconductores han revolucionado el mundo de la electrónica y han permitido grandes avances
en casi todas las ciencias; el más simple de los semiconductores es el diodo. El diodo solo
permite la circulación de la corriente en un único sentido, como se aprecia en la
figura~\ref{fig:diodo-func}, donde en la primera configuración permite el paso de corriente y en
la segunda la evita, principio del rectificador.

\begin{figure}[H]
    \centering
    \includegraphics[width=0.75\textwidth]{image03.png}
    \caption{Funcionamiento del diodo.}
    \label{fig:diodo-func}
\end{figure}

Los semiconductores son elementos que en condiciones normales son aislantes, pero con cierta
modificación en su organización molecular pueden convertirse en conductores. Esta modificación
consiste básicamente en el dopado con otros elementos que poseen enlaces covalentes incompletos;
ejemplos de semiconductores son el silicio, el selenio y el germanio. El símbolo del diodo y el
dispositivo se pueden apreciar en la figura~\ref{fig:diodo-simb}.

\begin{figure}[H]
    \centering
    \includegraphics[width=0.46\textwidth]{image04.png}
    \caption{Imagen y símbolo del diodo.}
    \label{fig:diodo-simb}
\end{figure}

El diodo está compuesto por la unión de un cristal tipo P con otro tipo N, en donde las cargas
opuestas de la unión AB crean una pequeña barrera de oposición, que con una mínima tensión de
polarización entre \qtyrange{0.2}{0.7}{\volt} logra la conducción de corriente en el diodo.

\begin{figure}[H]
    \centering
    \includegraphics[width=0.95\textwidth]{image05.png}
    \caption{Esquema eléctrico del funcionamiento del diodo.}
    \label{fig:diodo-pn}
\end{figure}

\FloatBarrier
\subsubsection{El condensador o capacitor}

El condensador es un dispositivo capaz de almacenar pequeñas cantidades de energía, para
devolverlas al circuito cuando este lo requiera; esto se puede apreciar en la
figura~\ref{fig:cap-carga}, donde se ha cargado un condensador para luego, con dicha carga, poder
encender un LED. Los condensadores se han utilizado en diferentes aplicaciones que van desde
temporizadores electrónicos, circuitos de rectificación, fenómenos de resonancia, supresión de
ruido y corrección del factor de potencia, entre otras.

\begin{figure}[H]
    \centering
    \includegraphics[width=0.80\textwidth]{image06.png}
    \caption{Carga y descarga de un condensador a través de un diodo LED.}
    \label{fig:cap-carga}
\end{figure}

Físicamente un condensador está compuesto por un par de placas metálicas separadas por un
material aislante, denominado dieléctrico, que se puede apreciar en la
figura~\ref{fig:cap-partes}. El dieléctrico puede ser aire, papel, cerámica, mica, plásticos y
otros.

\begin{figure}[H]
    \centering
    \includegraphics[width=0.78\textwidth]{image07.png}
    \caption{Partes de un condensador.}
    \label{fig:cap-partes}
\end{figure}

La capacidad de un condensador es la propiedad de almacenar mayor o menor cantidad de
electricidad o carga, dada en la siguiente ecuación, donde $Q$ es la cantidad de carga
almacenada, $C$ es la capacidad del condensador con unidad faradios y $V$ es la tensión entre las
armaduras. El faradio es una unidad muy grande, por lo que se utilizan submúltiplos de este, que
se pueden apreciar en la figura~\ref{fig:cap-ecuacion}.

\begin{figure}[H]
    \centering
    \includegraphics[width=0.82\textwidth]{image08.png}
    \caption{Ecuación del condensador y submúltiplos más utilizados de los faradios.}
    \label{fig:cap-ecuacion}
\end{figure}

\begin{figure}[H]
    \centering
    \includegraphics[width=0.72\textwidth]{image09.png}
    \caption{Distintos tipos de condensadores.}
    \label{fig:cap-tipos}
\end{figure}

Existe una gran variedad de condensadores, que se diferencian por su capacidad, la frecuencia de
funcionamiento y la cantidad de voltaje con la que operan. Entre ellos tenemos los condensadores
de papel, los condensadores de plástico o poliéster, condensadores cerámicos, condensadores de
mica, condensadores electrolíticos y los condensadores variables; esto se puede apreciar en la
figura~\ref{fig:cap-tipos}, además de los símbolos circuitales de los condensadores.

Por lo general los valores de la capacidad y tensión de trabajo aparecen inscritos en la
superficie del condensador electrolítico; esto se puede apreciar en la
figura~\ref{fig:cap-electro}. Los condensadores cerámicos por lo general usan la misma forma de
lectura de la resistencia, como se aprecia en la figura~\ref{fig:cap-ceramico}.

\begin{figure}[H]
    \centering
    \includegraphics[width=0.70\textwidth]{image10.png}
    \caption{Lectura de capacitor electrolítico.}
    \label{fig:cap-electro}
\end{figure}

% En el .doc original la imagen y la tabla de codigos de tension comparten un
% unico pie de figura, por eso van juntas dentro del mismo entorno figure.
\begin{figure}[H]
    \centering
    \includegraphics[width=0.90\textwidth]{image11.png}\\[10pt]
    \small
    \renewcommand{\arraystretch}{1.25}
    \begin{tabular}{|>{\centering\arraybackslash}m{3.1cm}|>{\centering\arraybackslash}m{3.1cm}
                    |>{\centering\arraybackslash}m{3.1cm}|}
    \hline
    0G = \qty{4}{\volt} DC   & 0L = \qty{5.5}{\volt} DC & 0J = \qty{6.3}{\volt} DC\\
    \hline
    1A = \qty{10}{\volt} DC  & 1C = \qty{16}{\volt} DC  & 1E = \qty{25}{\volt} DC\\
    \hline
    1H = \qty{50}{\volt} DC  & 1J = \qty{63}{\volt} DC  & 1K = \qty{80}{\volt} DC\\
    \hline
    2A = \qty{100}{\volt} DC & 2Q = \qty{110}{\volt} DC & 2B = \qty{125}{\volt} DC\\
    \hline
    2C = \qty{160}{\volt} DC & 2Z = \qty{180}{\volt} DC & 2D = \qty{200}{\volt} DC\\
    \hline
    2P = \qty{220}{\volt} DC & 2E = \qty{250}{\volt} DC & 2F = \qty{315}{\volt} DC\\
    \hline
    2V = \qty{350}{\volt} DC & 2G = \qty{400}{\volt} DC & 2W = \qty{450}{\volt} DC\\
    \hline
    2H = \qty{500}{\volt} DC & 2J = \qty{630}{\volt} DC & 3A = \qty{1000}{\volt} DC\\
    \hline
    \end{tabular}
    \renewcommand{\arraystretch}{1}
    \caption{Lectura de capacitores cerámicos.}
    \label{fig:cap-ceramico}
\end{figure}

\FloatBarrier
%=======================================================================
\section{Cuestionario previo}
%=======================================================================
\begin{enumerate}
\item Al realizar una medida de tensión alterna con un osciloscopio, aparece la imagen a
continuación.

% El .doc no da pie de figura a esta imagen, por eso va suelta y sin \caption.
\begin{center}
    \includegraphics[width=0.69\textwidth]{image12.png}
\end{center}

Complete la tabla a continuación:

\begin{table}[H]
    \centering
    \renewcommand{\arraystretch}{1.45}
    \begin{tabular}{|>{\centering\arraybackslash}m{6.0cm}|>{\centering\arraybackslash}m{4.0cm}|}
    \hline
    \textbf{Magnitud} & \textbf{Valor}\\
    \hline
    $V$ máximo ($V_{\max}$) & \\
    \hline
    $V$ eficaz ($V_{e}=V_{\max}/\sqrt{2}$) & \\
    \hline
    Periodo ($T$) & \\
    \hline
    Frecuencia ($f=1/T$) & \\
    \hline
    \end{tabular}
\end{table}

\item Complete la siguiente tabla de lectura de condensadores cerámicos.

\begin{table}[H]
    \centering
    \renewcommand{\arraystretch}{1.6}
    \begin{tabular}{|>{\centering\arraybackslash}m{1.8cm}|>{\centering\arraybackslash}m{2.2cm}
                    |>{\centering\arraybackslash}m{3.6cm}|>{\centering\arraybackslash}m{2.4cm}
                    |>{\centering\arraybackslash}m{2.2cm}|}
    \hline
    \textbf{Ejercicio} & \textbf{Código} & \textbf{Valor del condensador} &
    \textbf{Tolerancia} & \textbf{Voltaje}\\
    \hline
    1 & 2A 102J & & & \\
    \hline
    2 & 1H 103H & & & \\
    \hline
    3 & 2J 472K & & & \\
    \hline
    4 & 2D 104M & & & \\
    \hline
    \end{tabular}
\end{table}

\item Describa con sus propias palabras cómo funciona el dispositivo regulador LM317.

\lineasresp{2}
\end{enumerate}

\FloatBarrier
%=======================================================================
\section{Equipos y materiales}
%=======================================================================
\noindent\textbf{Laboratorio:} Electrónica y Comunicaciones.

\vspace{0.25cm}
\noindent\textbf{Equipos y dispositivos:}
\begin{itemize}
    \item 1 computador personal.
    \item 1 EEB.
    \item 1 diodo LED.
    \item 1 resistencia de \qty{1}{\kilo\ohm}.
    \item 1 potenciómetro.
    \item 1 bocina.
    \item 4 diodos.
    \item 1 capacitor.
\end{itemize}

\noindent\textbf{Software:}
\begin{itemize}
    \item WaveForms 2015.
\end{itemize}

\FloatBarrier
%=======================================================================
\section{Actividades}
%=======================================================================
\subsection{Actividad: conectando el osciloscopio y el generador de señales}

\begin{enumerate}[label=\textbf{\Alph*.}, leftmargin=*, labelsep=0.6em, series=actosc]
\item Arme el circuito de la figura~\ref{fig:circ1} en el EE Board y corra el programa WaveForms;
antes de iniciarlo, conecte la EE Board al ordenador, conecte su transformador y encienda el
\emph{switch} (ON).
\end{enumerate}

\begin{figure}[H]
    \centering
    \includegraphics[width=0.67\textwidth]{image13.png}
    \caption{Circuito para el desarrollo de la primera parte.}
    \label{fig:circ1}
\end{figure}

\begin{enumerate}[resume*=actosc]
\item Primero seleccione la opción \emph{Scope}, ver figura~\ref{fig:scope-sel}, y luego solo
habilite el canal 4 del osciloscopio, con los valores configurados en la
figura~\ref{fig:scope-cfg}.
\end{enumerate}

\begin{figure}[H]
    \centering
    \includegraphics[width=0.95\textwidth]{image14.png}
    \caption{Selección del osciloscopio y generador de señales para esta actividad.}
    \label{fig:scope-sel}
\end{figure}

\begin{figure}[H]
    \centering
    \includegraphics[width=0.93\textwidth]{image15.png}
    \caption{Configuración inicial del osciloscopio.}
    \label{fig:scope-cfg}
\end{figure}

\begin{enumerate}[resume*=actosc]
\item Luego seleccione la opción \emph{Wavegen}, ver figura~\ref{fig:scope-sel}, y configure el
generador con los valores que se muestran en la figura~\ref{fig:gen-cfg}; después encienda el
generador con el botón \emph{Run}.
\end{enumerate}

\begin{figure}[H]
    \centering
    \includegraphics[width=0.95\textwidth]{image16.png}
    \caption{Configuración inicial del generador de señales.}
    \label{fig:gen-cfg}
\end{figure}

\begin{enumerate}[resume*=actosc]
\item Conteste: ¿qué ocurre en el diodo LED?

\item Seleccione la opción \emph{Scope 1}, conecte la punta del osciloscopio en el punto central
entre el diodo y la resistencia, como se muestra en la figura~\ref{fig:scope-con}, y configure
los parámetros del osciloscopio como en la figura~\ref{fig:scope-align}.
\end{enumerate}

\begin{figure}[H]
    \centering
    \includegraphics[width=0.60\textwidth]{image17.png}
    \caption{Conexión del osciloscopio.}
    \label{fig:scope-con}
\end{figure}

\begin{figure}[H]
    \centering
    \includegraphics[width=0.95\textwidth]{image18.png}
    \caption{Configuración y alineación temporal del osciloscopio.}
    \label{fig:scope-align}
\end{figure}

\noindent\textbf{Preguntas:}
\begin{enumerate}
    \item ¿Qué ha pasado con la señal sinusoidal? Explique por qué se comporta de esa forma.
    \item Cambie la señal sinusoidal por una señal cuadrada, triangular, rampa y ruido. ¿Qué
          ocurre con la respuesta del circuito?
          \begin{enumerate}
              \item Para mayor claridad, conecte el osciloscopio directamente a la salida del
                    generador de señales, para ver las diferentes formas de onda.
          \end{enumerate}
    \item En el circuito de la figura~\ref{fig:circ1}, invierta la posición del diodo LED e
          inmediatamente cambie a un voltaje negativo en la amplitud del generador de señales.
          ¿Qué ha ocurrido en el circuito? ¿Por qué enciende el LED?
    \item Cambie el diodo LED por la bocina y establezca una frecuencia de \qty{2}{\kilo\hertz}.
          ¿Qué ocurre?
    \item Cambie los valores de la frecuencia hasta que deje de escuchar. ¿Cuál es el valor de
          frecuencia máximo que puede escuchar?
\end{enumerate}

\FloatBarrier
\subsection{Actividad: el rectificador}

\begin{enumerate}[label=\textbf{\Alph*.}, leftmargin=*, labelsep=0.6em, series=actrec]
\item En este primer apartado arme en el EE Board el circuito de la figura~\ref{fig:rect1} y
coloque el osciloscopio en el capacitor de \qty{10}{\micro\farad}. ¿Qué ocurre con la señal
sinusoidal de la entrada?
\end{enumerate}

\begin{figure}[H]
    \centering
    \includegraphics[width=0.95\textwidth]{image19.png}
    \caption{Esquemático del circuito a implementar.}
    \label{fig:rect1}
\end{figure}

\begin{enumerate}[resume*=actrec]
\item En este segundo apartado arme en el EE Board el circuito de la figura~\ref{fig:rect2} y
coloque el osciloscopio en el capacitor. ¿Qué ocurre con la señal sinusoidal de la entrada? ¿Qué
diferencia existe con la respuesta del apartado anterior?
\end{enumerate}

\begin{figure}[H]
    \centering
    \includegraphics[width=0.79\textwidth]{image20.png}
    \caption{Esquemático del circuito a implementar.}
    \label{fig:rect2}
\end{figure}

\FloatBarrier
%=======================================================================
\section{Tarea asignada}
%=======================================================================
\begin{itemize}
    \item Ninguna, es un taller.
\end{itemize}

\end{document}
